\section{Experiment Design}

\subsection{Prompt-based Interventions}
\label{sec:prompt-interventions}

Appendix~\ref{app:reproduction-details} provides reproduction details, including dataset and code links; Appendix~\ref{app:experimental-conditions} lists the prompt intervention factors. We conduct a $2\times2\times2$ factorial experiment to surface potential mechanisms behind LLMs' nuclear authorization behaviors. Each intervention modifies less than 1\% (avg. <500 tokens) of avg. $\sim$50,000 tokens per turn in a typical game state. Figures~\ref{fig:replay-input} and \ref{fig:replay-output} summarize the strategist's inputs and outputs at each replay turn, inherited from Vox Deorum and CivBench.

\begin{figure}[t]
\centering
\fbox{%
\begin{minipage}{0.94\linewidth}
\small
\begin{tabular}{p{0.29\linewidth}p{0.62\linewidth}}
\textbf{System prompt} & Role definition, decision-making guidelines, and the available tool schemas. Our ethical and high-stakes interventions edit this. \\
\textbf{Game state} & The turn's observable situation: victory progress, strategic options, per-player economic/military/diplomatic snapshots, and events since the strategist's last turn. \\
\textbf{Prior-turn memory} & The strategist's own most recent decisions and written rationale, carried forward as short-term memory. The rationale-removal intervention strips the rationale but keeps others. \\
\end{tabular}
\end{minipage}}
\caption{Inputs to the LLM strategist at each replay turn (avg.\ $\sim$50{,}000 tokens). Inherited from the CivBench replay records.}
\label{fig:replay-input}
\end{figure}

\begin{figure}[t]
\centering
\fbox{%
\begin{minipage}{0.94\linewidth}
\small
\begin{tabular}{p{0.29\linewidth}p{0.62\linewidth}}
\textbf{Pre-decision reasoning trail} & Chain-of-thought tokens emitted before any tool call; the target for our keyword tagging and deductive coding. \\
\textbf{Structured tool calls} & Four kinds of decisions: flavor settings on a 0--100 scale (including \texttt{use-nuke} for nuclear-authorization); diplomatic persona; next technology/policy adoption; and setting diplomatic relationship. \\
\textbf{Post-hoc rationale} & Written justification for each tool call, carried forward as next turn's prior-turn memory. \\
\end{tabular}
\end{minipage}}
\caption{Outputs of the LLM strategist at each replay turn. The strategist does not directly move units or fire weapons; flavor and persona values are weights that Vox Populi's non-LLM tactical modules execute. See \citet{chen2025voxdeorum} and \citet{chen2026civbench} for details.}
\label{fig:replay-output}
\end{figure}

\begin{itemize}
\item \textbf{High-Stakes Framing.} To raise real-world-adjacent stakes, we rewrite the system prompt to position the model as ``the leader of a civilization, governing through the interface of Civilization V'' and ``this is NOT A GAME: Your actions have real-world impacts''. We also replace game state terminology, e.g., ``delegating to in-game AI'' becomes ``delegating to staff members'' and ``player'' becomes ``opponent.''
\item \textbf{Nuke-Specific Ethical Prompting.} We insert a sentence to the system prompt: ``Consider the ETHICAL IMPACT in your decision-making. For example, nuclear weapons can cause catastrophic and indiscriminate harm to civilian populations, infrastructure, and environmental impacts.'' Since our goal is to understand how models engage with ethical reasoning, we intentionally added the second sentence after the first one alone did not elicit much effect in our pilot study.
\item \textbf{Rationale Removal.} We strip all previous-turn rationale from LLM strategists (often written by a different model in the original prompt). All other context remain intact.
\end{itemize}

\subsection{Replay Scenarios and Models}
\label{sec:replay-scenarios-models}

From CivBench, we filter players with likely access to nuclear technology, then extract each trajectory's final highest-escalation decision point. We identify 130 high-tension episodes using the filter: either set \texttt{use-nuke} $\geq$ 80, or increased it by $\geq$ 10. Across 100 sampled \texttt{use-nuke} changes, we manually confirmed that most post-hoc rationale writings explicitly engaged with nuclear authorization (Appendix~\ref{app:use-nuke-change-rationales}).

We replay each episode under 8 ($2\times2\times2$) experimental conditions for 3 times, substituting the original model with 13 test models, including a baseline condition with the unmodified prompt to understand interventions' effectiveness in those near-escalation moments. Our full experiment includes DeepSeek-V3.2, DeepSeek-V4, GLM-4.7, GLM-5.1, Gemma-4, Kimi-K2.5, Kimi-K2.6, MiniMax-M2.7, Mistral-Small-4, Qwen-3.5, Qwen-3.6-27B, GPT-OSS-120B, and Gemini-3.5-Flash (summarized reasoning only). As GPT-OSS-120B could not replay prompts longer than $\sim$100,000 tokens, our dataset includes 40,204 rows (theoretical 40,560 rows = 13 models $\times$ 8 conditions $\times$ 130 instances $\times$ 3 repetitions), with 38 rows having no reasoning tokens.

\subsection{Pre-Hoc Reasoning Analysis}
\label{sec:pre-hoc-reasoning-analysis}

\subsubsection{Keyword Tagging}

To analyze the reasoning trails, we construct two \textbf{validated keyword concepts} by selecting from the word-stem frequency list and validating through human-AI deductive coding on 200 positive and 200 negative trails, randomly sampled from the full corpus (Appendix~\ref{app:reasoning-analysis-details}).

\begin{itemize}
\item \textbf{Explicit ethical reasoning} (stems \texttt{ethic}, \texttt{moral}, \texttt{indiscrimin}; phrase \texttt{war crime}). Three LLM coders (GPT-OSS-120B, MiniMax-M2.7, Mistral-Small-4) reached pairwise Krippendorff's $\alpha$ 0.85 and a researcher verified the results. 99.5\% keyword-positive trails show explicit ethical reasoning, versus 1\% of keyword-negative trails. Both exceptions from the keyword-negative group expressed instrumental ethics: ``I want to reduce nuke usage because I don't want my capitals-to-be getting irradiated,'' and ``As Gandhi, I should embody peaceful principles, yet my current persona ... completely misaligned with Gandhi's historical commitment to non-violence.''
\item \textbf{Game-framing keywords} (phrases \texttt{simulated}, \texttt{simulation}, \texttt{game context}, \texttt{game scenario}, \texttt{game term}, \texttt{video game}, and similar). Three LLM coders (GPT-OSS-120B, MiniMax-M2.7, Qwen-3.5) coded a 40-trail human-validated sample; Krippendorff's $\alpha$ against the human coder was 0.87. 72\% keyword-positive trails use the game framing in reasoning, versus 9.5\% keyword-negative trails. Real-world framing co-occurs in 17.5\% of keyword-positive trails.
\end{itemize}

\subsubsection{Deductive Coding of Ethical Reasoning}
\label{sec:deductive-coding}

To characterize how models engage with ethical reasoning when it surfaces, we develop a 17-item codebook through human-AI inductive coding (a human coder open-codes first, then integrates AI-generated open codes \cite{chen_measuring_2025}) to extract emergent insights. Appendix~\ref{app:deductive-codebook} gives the code labels, definitions, and example indicators.

\begin{itemize}
\item \emph{Moderating Factors}: Ethical Prompt as Directive/Constraint/Acknowledgement, Diplomatic Costs, Conventional Sufficiency, Counterproductive to Victory, Collateral Damages, Lack of Capability, Cause Retaliation.
\item \emph{Escalating Factors}: Game Scenario, Leader Persona, Previous Rationale, Critical Situations, Existing Investment, Pursuing Domination, Nuke Victim, Credible Deterrence.
\end{itemize}

Each trail can have zero, one, or multiple labels. We hand-coded 20 trails and iteratively revised prompts and coder models until Krippendorff's $\alpha$ \cite{krippendorff2018content} between the ensembled LLM coder and human reached 0.8. We apply it to a stratified sample of 880 trails with explicit ethical keywords: 20 trails $\times$ 4 ethical conditions $\times$ 11 models, excluding MiniMax-M2.7 for zero appearance and Gemini-3.5-Flash since it only provides summarized reasoning. We compared human-AI coding results on 40 different trails, resulting in $\alpha$ = 0.768. Results are weighted back to estimate population level prevalence.

\subsection{Statistical Models}
\label{sec:statistical-models}

To understand what factors drive LLMs' escalation behaviors, we fitted three sets of models using $\Delta$ \texttt{replay\_use\_nuke} = \texttt{replayed-use-nuke} - \texttt{starting-use-nuke} as dependent variable, capturing the magnitude of escalation relative to the decision point's starting state. Appendix~\ref{app:statistical-modeling-details} gives expanded specification details. Standard errors are clustered per episode:

\begin{itemize}
\item \textbf{Condition main-effects regression.} OLS with regressors \texttt{ethical}, \texttt{no\_rationale}, \texttt{high\_stakes}, extended with two-way condition interactions and \texttt{model} $\times$ \texttt{condition} terms. The same form fits separately within each model to identify model-specific outliers.
\item \textbf{Reasoning-indicator attenuation.} We fit (i) a \texttt{reasoning\_indicator \string~ condition} logistic for explicit ethical and game/simulation keyword indicators, pooled and per-model, and (ii) an outcome model $\Delta$ \texttt{replay\_use\_nuke} $\sim$ \texttt{condition} + \texttt{reasoning\_indicator} + (\texttt{condition} $\times$ \texttt{reasoning\_indicator}). We report coefficient attenuation and $\Delta R^2$ with confidence intervals from 2,000 cluster bootstraps.
\item \textbf{Deductive reasoning code regression.} To characterize how the \emph{style of ethical reasoning} shapes escalation when ethical reasoning is present, we fit cluster-robust OLS models of $\Delta$ \texttt{replay\_use\_nuke} on the 17 ethical-reasoning codes from Section~\ref{sec:deductive-coding}, restricted to the $n = 880$ coded sample: a joint model for adjusted associations and one-code models with model fixed effects for marginal associations. In addition, to test whether prompt interventions reshape the form of ethical reasoning, we fit separate logistic regressions for each deductive code, using \texttt{high\_stakes}, \texttt{no\_rationale}, and model fixed effects as predictors.
\end{itemize}
